\chapter{H\"older Estimates for the Gradient}

Let $\Omega\subset\R^n$ be bounded and consider
\begin{equation}
Qu=a^{ij}(x,u,Du)D_{ij}u+b(x,u,Du)=0.
\tag{13.1}
\end{equation}

\section{Equations of Divergence Form}

Suppose first that $Q$ is equivalent to
\begin{equation}
Qu=\operatorname{div}A(x,u,Du)+B(x,u,Du)=0,
\tag{13.2}
\end{equation}
where $A\in C^1$ and $B\in C^0$.  A weak solution satisfies
\[
\int_\Omega\{A(x,u,Du)\cdot D\zeta-B(x,u,Du)\zeta\}\,dx=0
\quad(\zeta\in C_0^1(\Omega)).
\]
For
\[
\delta_k A^i=D_{x_k}A^i+p_kD_zA^i,
\quad
\bar a^{ij}=D_{p_j}A^i(x,u,Du),
\quad
f_k^i=\delta_kA^i(x,u,Du)+\delta_k^iB(x,u,Du),
\]
the derivative $w=D_ku$ satisfies
\begin{equation}
\int_\Omega(\bar a^{ij}D_jw+f_k^i)D_i\zeta\,dx=0,
\quad \zeta\in C_0^1(\Omega).
\tag{13.3}
\end{equation}
For $|z|+|p|\le K$, choose constants such that
\begin{equation}
0<\lambda_K\le\lambda(x,z,p),\qquad
\Lambda_K\ge |D_{p_j}A^i(x,z,p)|,
\qquad
\mu_K\ge |\delta_jA^i(x,z,p)|+|B(x,z,p)|.
\tag{13.4}
\end{equation}

\begin{theorem}[Interior divergence-form gradient estimate]
\label{thm:gt13-holder-gradient-13-1}\dcref{ch13:13.1}\group{ch13-holder-gradient}
\source{gilbarg-trudinger:page-0332}
Let $u\in C^2(\Omega)$ solve \textup{(13.2)}, and assume that $Q$ is elliptic,
$A\in C^1(\Omega\times\R\times\R^n)$, and
$B\in C^0(\Omega\times\R\times\R^n)$.  If $\Omega'\Subset\Omega$, then
\begin{equation}
[Du]_{\alpha;\Omega'}\le Cd^{-\alpha},
\qquad d=\operatorname{dist}(\Omega',\partial\Omega),
\tag{13.5}
\end{equation}
where $K=|u|_{1;\Omega}$,
\[
C=C(n,K,\Lambda_K/\lambda_K,\mu_K/\lambda_K,
\operatorname{diam}\Omega),
\qquad
\alpha=\alpha(n,\Lambda_K/\lambda_K)>0.
\]
\end{theorem}
\begin{proof}
\sketch Equation \textup{(13.3)} is a uniformly elliptic divergence-form
equation for each $D_ku$, with coefficient and inhomogeneous-term bounds
given by \textup{(13.4)}.  The interior oscillation estimate for such equations
applied on balls of radius comparable to $d$ gives \textup{(13.5)}.
\end{proof}

Near $x_0\in\partial\Omega$, flatten the boundary by $y=\psi(x)$ and set
$v=u\circ\psi^{-1}$.  After subtracting the boundary datum, $v=0$ on the
flat boundary and
\begin{equation}
\bar Qv=D_{y_i}\bar A^i+\bar B=0,
\quad
\bar A^i=\frac{\partial y_i}{\partial x_r}A^r,
\quad
\bar B=-D_{y_i}\!\left(\frac{\partial y_i}{\partial x_r}\right)A^r+B.
\tag{13.6}
\end{equation}
The boundary linear estimate gives, for tangential derivatives,
\begin{equation}
[D_{y_k}v]_{\alpha;D'\cap D^+}\le C,
\qquad k<n.
\tag{13.7}
\end{equation}
A Caccioppoli estimate then yields
\begin{equation}
\int_{B_R}|D(D_{y_k}v)|^2\,dy\le CR^{n-2+2\alpha},
\qquad k<n.
\tag{13.8}
\end{equation}
Solving the transformed equation for $D_{y_ny_n}v$ extends this estimate to
the normal derivative, and Morrey's estimate gives
\begin{equation}
[Du]_{\alpha;B'\cap\Omega}\le C.
\tag{13.9}
\end{equation}

\begin{theorem}[Global divergence-form gradient estimate]
\label{thm:gt13-holder-gradient-13-2}\dcref{ch13:13.2}\group{ch13-holder-gradient}
\source{gilbarg-trudinger:page-0335}
Assume that $u\in C^2(\overline\Omega)$ solves \textup{(13.2)}, that $Q$ is
elliptic on $\overline\Omega$, that $A\in C^1$ and $B\in C^0$ on
$\overline\Omega\times\R\times\R^n$, and that $\partial\Omega\in C^2$.
If $u=\varphi$ on $\partial\Omega$ with
$\varphi\in C^2(\overline\Omega)$, then
\begin{equation}
[Du]_{\alpha;\Omega}\le C,
\tag{13.10}
\end{equation}
where $K=|u|_{1;\Omega}$, $\Phi=|\varphi|_{2;\Omega}$,
\[
C=C(n,K,\Lambda_K/\lambda_K,\mu_K/\lambda_K,\Omega,\Phi),
\qquad
\alpha=\alpha(n,\Lambda_K/\lambda_K,\Omega)>0.
\]
\end{theorem}
\begin{proof}
\sketch Apply Theorem~\ref{thm:gt13-holder-gradient-13-1} away from the
boundary and \textup{(13.6)--(13.9)} in finitely many boundary charts.
\end{proof}

The estimates \textup{(13.5)} and \textup{(13.10)} remain valid for
$u\in C^{0,1}(\overline\Omega)\cap W^{2,2}(\Omega)$ and
$\varphi\in W^{2,q}(\Omega)$, $q>n$, with
$\Phi=\|\varphi\|_{W^{2,q}(\Omega)}$ and with $\alpha$ also depending on $q$.

\section{Equations in Two Variables}

For a solution of \textup{(13.1)} in $\Omega\subset\R^2$, the derivatives
$w_1=D_1u$ and $w_2=D_2u$ satisfy
\begin{equation}
\begin{aligned}
D_1\!\left(\frac{a^{11}}{a^{22}}D_1w_1
 +\frac{2a^{12}}{a^{22}}D_2w_2\right)+D_{22}w_1
 &=-D_1\frac b{a^{22}},\\
D_{11}w_2+D_2\!\left(\frac{2a^{12}}{a^{11}}D_1w_2
 +\frac{a^{22}}{a^{11}}D_2w_2\right)
 &=-D_2\frac b{a^{11}}.
\end{aligned}
\tag{13.11}
\end{equation}
For $|z|+|p|\le K$, take
\begin{equation}
0<\lambda_K\le\lambda(x,z,p),\qquad
\Lambda_K\ge|a^{ij}(x,z,p)|,\qquad
\mu_K\ge|b(x,z,p)|.
\tag{13.12}
\end{equation}

\begin{theorem}[Two-dimensional interior estimate]
\label{thm:gt13-holder-gradient-13-3}\dcref{ch13:13.3}\group{ch13-holder-gradient}
\source{gilbarg-trudinger:page-0335}
Let $u\in C^2(\Omega)$ solve \textup{(13.1)} in $\Omega\subset\R^2$.
Assume $Q$ is elliptic and $a^{ij},b$ are continuous on
$\Omega\times\R\times\R^2$.  For every $\Omega'\Subset\Omega$,
\begin{equation}
[Du]_{\alpha;\Omega'}\le Cd^{-\alpha},
\qquad d=\operatorname{dist}(\Omega',\partial\Omega),
\tag{13.13}
\end{equation}
where
\[
C=C(K,\Lambda_K/\lambda_K,\mu_K/\lambda_K,
\operatorname{diam}\Omega),
\qquad
\alpha=\alpha(\Lambda_K/\lambda_K)>0.
\]
\end{theorem}
\begin{proof}
\sketch Apply the divergence-form interior estimate to the two equations
in \textup{(13.11)}.
\end{proof}

\begin{theorem}[Two-dimensional global estimate]
\label{thm:gt13-holder-gradient-13-4}\dcref{ch13:13.4}\group{ch13-holder-gradient}
\source{gilbarg-trudinger:page-0336}
Let $u\in C^2(\overline\Omega)$ solve \textup{(13.1)} in
$\Omega\subset\R^2$.  Assume $Q$ is elliptic on $\overline\Omega$, the
coefficients $a^{ij},b$ are continuous on
$\overline\Omega\times\R\times\R^2$, $\partial\Omega\in C^2$, and
$u=\varphi$ on $\partial\Omega$ with
$\varphi\in C^2(\overline\Omega)$.  Then
\begin{equation}
[Du]_{\alpha;\Omega}\le C,
\tag{13.14}
\end{equation}
where $K=|u|_{1;\Omega}$, $\Phi=|\varphi|_{2;\Omega}$,
\[
C=C(K,\Lambda_K/\lambda_K,\mu_K/\lambda_K,\Omega,\Phi),
\qquad
\alpha=\alpha(\Lambda_K/\lambda_K,\Omega)>0.
\]
\end{theorem}
\begin{proof}
\sketch Combine \textup{(13.11)} with the global divergence-form estimate.
The weak-solution extension following Theorem~\ref{thm:gt13-holder-gradient-13-2}
applies as well.
\end{proof}

\section{Equations of General Form: Interior Estimate}

Assume $a^{ij}\in C^1$, $b\in C^0$, and initially $u\in C^3$.  With
$\delta_kg=D_{x_k}g+p_kD_zg$, differentiation of \textup{(13.1)} gives
\begin{equation}
a^{ij}D_{kij}u+D_{p_l}a^{ij}D_{lk}uD_{ij}u
+\delta_ka^{ij}D_{ij}u+D_kb=0,
\tag{13.15}
\end{equation}
or, in divergence form,
\begin{equation}
D_i(a^{ij}D_{kj}u)
+(D_{p_l}a^{ij}-D_{p_j}a^{il})D_{lk}uD_{ij}u
+\delta_ka^{ij}D_{ij}u-\delta_ia^{ij}D_{kj}u+D_kb=0.
\tag{13.16}
\end{equation}
For $v=|Du|^2$, multiplication by $D_ku$ and summation yield
\begin{equation}
\begin{aligned}
0={}&-a^{ij}D_{ki}uD_{kj}u+\tfrac12D_i(a^{ij}D_jv)
+\tfrac12(D_{p_l}a^{ij}-D_{p_j}a^{il})D_{ij}uD_lv\\
&+D_ku\,\delta_ka^{ij}D_{ij}u
-\tfrac12\delta_ia^{ij}D_jv+D_k(bD_ku)-b\Delta u.
\end{aligned}
\tag{13.17}
\end{equation}
Set
\begin{equation}
w=w_r=\gamma D_ru+v.
\tag{13.18}
\end{equation}
Combining \textup{(13.16)} and \textup{(13.17)} gives
\begin{equation}
\begin{aligned}
0={}&-2a^{ij}D_{ki}uD_{kj}u
+D_i\bigl(a^{ij}D_jw+(2D_iu+\gamma\delta^{ir})b\bigr)\\
&+(D_{p_l}a^{ij}-D_{p_j}a^{il})D_{ij}uD_lw
+\bigl[(2D_ku+\gamma\delta^{kr})\delta_ka^{ij}-2b\delta^{ij}\bigr]D_{ij}u
-\delta_ia^{ij}D_jw.
\end{aligned}
\tag{13.19}
\end{equation}
Its weak form is
\begin{equation}
\int_\Omega\!\left\{(\bar a^{ij}D_jw+f^i)D_i\zeta
+(2\bar a^{ij}D_{ki}uD_{kj}u-a_l^{\,ij}D_{ij}uD_lw
+c^iD_iw-b^{ij}D_{ij}u)\zeta\right\}dx=0.
\tag{13.20}
\end{equation}
This identity extends by $C^2$ approximation, and also to
$C^{0,1}\cap W^{2,2}$ solutions.  Ellipticity and Schwarz's inequality imply
\begin{equation}
\int_\Omega(\bar a^{ij}D_jw+f^i)D_i\zeta\,dx
\le\int_\Omega\left\{
\left(\lambda+\lambda^{-1}\sum|a^{ij}|^2\right)|Dw|^2
+\lambda^{-1}\sum(|c^i|^2+|b^{ij}|^2)\right\}\zeta\,dx
\tag{13.21}
\end{equation}
for nonnegative $\zeta$.  Testing with $e^{2\chi w}\zeta$, where
$\chi=\sup(1+\lambda^{-2}\sum|a^{ij}|^2)$, gives
\begin{equation}
\int_\Omega\{(\widetilde a^{ij}D_iw+\widetilde f^j)D_j\zeta
-\widetilde g\zeta\}\,dx\le0,
\tag{13.22}
\end{equation}
equivalently
\begin{equation}
D_i(\widetilde a^{ij}D_jw)\ge-(\widetilde g+D_i\widetilde f^i)
\tag{13.23}
\end{equation}
in the generalized sense.

Let $M=\sup|Du|$, choose $\gamma=10nM$, and put
$w_r^\pm=\pm\gamma D_ru+v$.  If $D_ru$ has maximal oscillation among the
components of $Du$ on a set $\mathcal G$, then
\begin{equation}
8nM\,\operatorname{osc}_{\mathcal G}D_ru
\le\operatorname{osc}_{\mathcal G}w_r^\pm
\le12nM\,\operatorname{osc}_{\mathcal G}D_ru.
\tag{13.24}
\end{equation}
Writing $W^\pm=\sup_{\mathcal G}w_r^\pm$, one also has
\begin{equation}
\inf_{\mathcal G}\sum_\pm(W^\pm-w_r^\pm)
\ge6nM\,\operatorname{osc}_{\mathcal G}D_ru
\ge\tfrac12\operatorname{osc}_{\mathcal G}w_r^\pm.
\tag{13.25}
\end{equation}
For $|z|+|p|\le K$, choose
\begin{equation}
0<\lambda_K\le\lambda(x,z,p),\qquad
\mu_K\ge |a^{ij}|+|D_{p_k}a^{ij}|+|D_za^{ij}|+|D_{x_k}a^{ij}|+|b|.
\tag{13.26}
\end{equation}
The weak Harnack inequality on $B_{4R}$ gives
\begin{equation}
R^{-n}\int_{B_{2R}}(W^\pm-w_r^\pm)\,dx
\le C\{W^\pm-\sup_{B_R}w_r^\pm+\sigma(R)\},
\qquad \sigma(R)=R+R^2,
\tag{13.27}
\end{equation}
and hence an oscillation recurrence of the form used below.

\begin{lemma}[Oscillation iteration]
\label{lem:gt13-holder-gradient-13-5}\dcref{ch13:13.5}\group{ch13-holder-gradient}
\source{gilbarg-trudinger:page-0340}
Let $\omega_i,\bar\omega_i\ge0$, $i=1,\ldots,N$, be nondecreasing on
$(0,R_0)$.  Suppose that for each $R\le R_0$ there is an index $r$ such that
\[
\bar\omega_r(R)\ge\delta_0\omega_i(R)\quad(i=1,\ldots,N)
\]
and
\begin{equation}
\bar\omega_r(\delta R)\le\gamma\bar\omega_r(R)+\sigma(R),
\tag{13.28}
\end{equation}
where $\sigma$ is nondecreasing, $\delta_0>0$, and $0<\delta,\gamma<1$.
Then, for every $\mu\in(0,1)$ and $R\le R_0$,
\begin{equation}
\omega_i(R)\le C\left\{
\left(\frac R{R_0}\right)^\alpha\omega_0
+\sigma(R^\mu R_0^{1-\mu})\right\},
\tag{13.29}
\end{equation}
where $\omega_0=\max_i\bar\omega_i(R_0)$,
$C=C(N,\delta_0,\delta,\gamma)$, and
$\alpha=\alpha(N,\delta_0,\delta,\gamma,\mu)>0$.
\end{lemma}
\begin{proof}
\sketch Iterate \textup{(13.28)} on the geometric radii
$R_0,\delta R_0,\ldots$ until the scale passes
$R^\mu R_0^{1-\mu}$.  Monotonicity controls the remaining scale change, and
the lower comparison with every $\omega_i$ gives \textup{(13.29)}.
\end{proof}

Applying Lemma~\ref{lem:gt13-holder-gradient-13-5} to the $2n$ functions
$w_r^\pm$, with $\delta=1/4$ and $\mu=1/2$, yields
\begin{equation}
\operatorname{osc}_{B_R(y)}D_iu\le CR^\alpha.
\tag{13.30}
\end{equation}

\begin{theorem}[Interior gradient estimate for general equations]
\label{thm:gt13-holder-gradient-13-6}\dcref{ch13:13.6}\group{ch13-holder-gradient}
\source{gilbarg-trudinger:page-0340}
Let $u\in C^2(\Omega)$ solve \textup{(13.1)}.  Assume $Q$ is elliptic,
$a^{ij}\in C^1(\Omega\times\R\times\R^n)$, and
$b\in C^0(\Omega\times\R\times\R^n)$.  For every
$\Omega'\Subset\Omega$,
\begin{equation}
[Du]_{\alpha;\Omega'}\le Cd^{-\alpha},
\tag{13.31}
\end{equation}
where $K=|u|_{1;\Omega}$, $d=\operatorname{dist}(\Omega',\partial\Omega)$,
\[
C=C(n,K,\mu_K/\lambda_K,\operatorname{diam}\Omega),
\qquad
\alpha=\alpha(n,K,\mu_K/\lambda_K)>0.
\]
\end{theorem}
\begin{proof}
\sketch The reduction \textup{(13.15)--(13.27)} and the iteration lemma
give \textup{(13.30)} on every interior ball.  A covering at scale $d$
gives \textup{(13.31)}.
\end{proof}

\section{Equations of General Form: Boundary Estimate}

A $C^2$ boundary-flattening map preserves the hypotheses, since
\begin{equation}
Qu=a^{kl}\frac{\partial y_i}{\partial x_k}
\frac{\partial y_j}{\partial x_l}D_{y_iy_j}u
+a^{ij}\frac{\partial^2y_k}{\partial x_i\partial x_j}D_{y_k}u+b.
\tag{13.32}
\end{equation}
After subtracting the boundary values, work near a flat boundary with
\begin{equation}
v'=\sum_{i=1}^{n-1}|D_iu|^2,
\qquad w=w_r=\gamma D_ru+v',\quad r<n.
\tag{13.33}
\end{equation}
The equation determines the missing second derivative by
\begin{equation}
D_{nn}u=-\frac1{a^{nn}}
\left(\sum_{(i,j)\ne(n,n)}a^{ij}D_{ij}u+b\right).
\tag{13.34}
\end{equation}
The boundary weak Harnack inequality gives
\begin{equation}
\operatorname{osc}_{D^+\cap B_R(y)}D_iu\le CR^\alpha,
\qquad i<n.
\tag{13.35}
\end{equation}
Testing the inequality for $w$ by $\eta^2(w-c)^+$ gives
\begin{equation}
\int_{B_R^+}|Dw|^2\,dx
\le CR^{n-2}\{R^2+(\operatorname{osc}_{B_{2R}}w)^2\}
\le CR^{n-2+2\alpha}.
\tag{13.36}
\end{equation}
Taking $\gamma=0,\pm1$ yields
\begin{equation}
\int_{D^+\cap B_R}|D(D_ru)|^2\,dx\le CR^{n-2+2\alpha},
\qquad r<n,
\tag{13.37}
\end{equation}
and consequently
\begin{equation}
\int_{D^+\cap B_R}|D_{ij}u|^2\,dx\le CR^{n-2+2\alpha}
\quad\text{when }j\ne n.
\tag{13.38}
\end{equation}
Equation \textup{(13.34)} supplies the remaining derivative, so Morrey's
estimate gives
\begin{equation}
[Du]_{\alpha;D'\cap D^+}\le C.
\tag{13.39}
\end{equation}

\begin{theorem}[Global gradient estimate for general equations]
\label{thm:gt13-holder-gradient-13-7}\dcref{ch13:13.7}\group{ch13-holder-gradient}
\source{gilbarg-trudinger:page-0343}
Let $u\in C^2(\overline\Omega)$ solve \textup{(13.1)}.  Assume $Q$ is
elliptic on $\overline\Omega$, $a^{ij}\in C^1$ and $b\in C^0$ on
$\overline\Omega\times\R\times\R^n$, $\partial\Omega\in C^2$, and
$u=\varphi$ on $\partial\Omega$ with
$\varphi\in C^2(\overline\Omega)$.  Then
\begin{equation}
[Du]_{\alpha;\Omega}\le C,
\tag{13.40}
\end{equation}
where $K=|u|_{1;\Omega}$, $\Phi=|\varphi|_{2;\Omega}$,
\[
C=C(n,K,\mu_K/\lambda_K,\Omega,\Phi),
\qquad
\alpha=\alpha(n,K,\mu_K/\lambda_K,\Omega)>0.
\]
\end{theorem}
\begin{proof}
\sketch Combine Theorem~\ref{thm:gt13-holder-gradient-13-6} on the interior
with \textup{(13.32)--(13.39)} in finitely many boundary charts.
\end{proof}

The estimates \textup{(13.31)} and \textup{(13.40)} remain valid for
$u\in C^{0,1}(\overline\Omega)\cap W^{2,2}(\Omega)$ and
$\varphi\in W^{2,q}(\Omega)$, $q>n$, with
$\Phi=\|\varphi\|_{W^{2,q}(\Omega)}$ and constants also depending on $q$.
If $T\subset\partial\Omega$ is a $C^2$ boundary portion, the hypotheses of
Theorem~\ref{thm:gt13-holder-gradient-13-7} hold on $\Omega\cup T$, and
$u=\varphi$ on $T$, then for $\Omega'\Subset\Omega\cup T$,
\begin{equation}
[Du]_{\alpha;\Omega'}\le C,
\tag{13.41}
\end{equation}
where $d=\operatorname{dist}(\Omega',\partial\Omega\setminus T)$,
\[
C=C(n,K,\mu_K/\lambda_K,T,\Phi,d),
\qquad
\alpha=\alpha(n,K,\mu_K/\lambda_K,T)>0.
\]

\section{Application to the Dirichlet Problem}

\begin{theorem}[Existence from a uniform gradient bound]
\label{thm:gt13-holder-gradient-13-8}\dcref{ch13:13.8}\group{ch13-holder-gradient}
\source{gilbarg-trudinger:page-0344}
Let $\Omega\subset\R^n$ be bounded.  Suppose that $Q$ is elliptic on
$\overline\Omega$, $a^{ij}\in C^1$ and $b\in C^\alpha$ on
$\overline\Omega\times\R\times\R^n$, $0<\alpha<1$, and that
$\partial\Omega\in C^{2,\alpha}$ and
$\varphi\in C^{2,\alpha}(\overline\Omega)$.  Assume there is an $M$,
independent of $u$ and $\sigma$, such that every
$u\in C^{2,\alpha}(\overline\Omega)$ solving
\begin{equation}
Q_\sigma u=a^{ij}(x,u,Du)D_{ij}u+\sigma b(x,u,Du)=0,
\quad u=\sigma\varphi\text{ on }\partial\Omega,\quad 0\le\sigma\le1,
\tag{13.42}
\end{equation}
satisfies
\begin{equation}
\|u\|_{C^1(\overline\Omega)}
=\sup_\Omega|u|+\sup_\Omega|Du|<M.
\tag{13.43}
\end{equation}
Then $Qu=0$ in $\Omega$, $u=\varphi$ on $\partial\Omega$, has a solution
in $C^{2,\alpha}(\overline\Omega)$.  If $Q$ is in divergence form or
$n=2$, it is enough to assume $a^{ij}\in C^\alpha$.

The same conclusion holds for a family
\begin{equation}
Q_\sigma u=a^{ij}(x,u,Du;\sigma)D_{ij}u+b(x,u,Du;\sigma)=0,
\quad u=\sigma\varphi\text{ on }\partial\Omega,\quad 0\le\sigma\le1,
\tag{13.44}
\end{equation}
provided $Q_1=Q$, $b(x,z,p;0)=0$, every $Q_\sigma$ is elliptic on
$\overline\Omega$, and the coefficients depend sufficiently smoothly on
$(x,z,p,\sigma)$.
\end{theorem}
\begin{proof}
\sketch The uniform $C^1$ bound and Theorems
\ref{thm:gt13-holder-gradient-13-2},
\ref{thm:gt13-holder-gradient-13-4}, and
\ref{thm:gt13-holder-gradient-13-7} give uniform H\"older control of $Du$.
Schauder estimates then bound the solution set in $C^{2,\alpha}$, and the
Leray--Schauder continuation theorem applies to \textup{(13.42)}.  The same
argument uses the parameterized continuation theorem for \textup{(13.44)}.
\end{proof}
